\documentclass{article}
\usepackage{fontspec}
\usepackage{listings}

\usepackage{orcidlink}
\usepackage{xcolor}
\usepackage{hyperref}

\hypersetup{
  colorlinks=true,
  linkcolor=blue,
  citecolor=blue,
  urlcolor=blue,
  pdftitle={Supplemental: Formalized Lean 4 Proofs for Contexts and Pseudocontexts},
  pdfauthor={Mirko Navara and Karl Svozil},
  pdfkeywords={pseudocontext, projection lattice, rank-one decomposition, Lean, formal verification}
}

% Use a monospaced font with broad Unicode coverage for Lean source.
\setmonofont{DejaVuSansMono.ttf}[
  Scale=MatchLowercase,
  BoldFont=DejaVuSansMono-Bold.ttf,
  ItalicFont=DejaVuSansMono-Oblique.ttf,
  BoldItalicFont=DejaVuSansMono-BoldOblique.ttf
]

\definecolor{keywordcolor}{rgb}{0.7, 0.1, 0.1}
\definecolor{commentcolor}{rgb}{0.2, 0.6, 0.2}
\definecolor{stringcolor}{rgb}{0.1, 0.1, 0.7}
\definecolor{codebackground}{rgb}{0.94, 0.94, 0.94}

\lstdefinelanguage{lean}{
  keywords={abbrev, axiom, by, calc, cases, constructor, def, end, exact,
    fun, funext, have, import, intro, namespace, open, rcases, rfl, rw,
    set_option, simp, theorem, variable, where},
  sensitive=true,
  comment=[l]{--},
  morecomment=[s]{/-}{-/},
  morestring=[b]",
}

\lstset{
  language=lean,
  basicstyle=\ttfamily\small,
  keywordstyle=\color{keywordcolor}\bfseries,
  commentstyle=\color{commentcolor}\itshape,
  stringstyle=\color{stringcolor},
  backgroundcolor=\color{codebackground},
  numbers=left,
  numberstyle=\tiny\color{gray},
  stepnumber=1,
  breaklines=true,
  showstringspaces=false,
  tabsize=2,
  frame=single,
  linewidth=\linewidth,
  xleftmargin=3pt,
  xrightmargin=3pt,
  framesep=4pt
}

\title{Supplemental: Formalized Lean 4 Proofs\\
Contexts and Pseudocontexts}
\author{Mirko Navara\,\orcidlink{0000-0002-0880-5992}
\and Karl Svozil\,\orcidlink{0000-0001-6554-2802}}
\date{June 2, 2026}

\begin{document}
\maketitle

\section*{Overview}
This document gives the complete Lean 4 source for the formal verification
accompanying the manuscript \emph{Contexts and pseudocontexts}.  The Lean
development formalizes the finite coordinate-operator core of the paper:
rank-one projectors, finite structure operators, the complex phase
construction, ray-disjointness, the revised genuineness certificate, and the
basic trace identities for rank-one decompositions.

The formalization is intentionally scoped to the algebraic spine of the
article.  It does not attempt to formalize the full projection-lattice
interface, Gleason-theorem infrastructure, Schur--Horn theory, or the
topological genericity arguments used elsewhere in the manuscript.  Instead,
it supplies a machine-checked certificate for the explicit finite-dimensional
operator calculations that are most directly amenable to Lean verification.

The source has been checked with Lean 4 \texttt{v4.30.0} and mathlib
\texttt{v4.30.0}.  In the archived workspace used for this supplement, the
verification command is
\[
  \texttt{powershell.exe -ExecutionPolicy Bypass -File .\textbackslash verify.ps1}.
\]
The Lean files contain no placeholders such as \texttt{sorry} or
\texttt{admit}.

The formalization consists of a single main Lean file:
\begin{itemize}
  \item Section~\ref{app:lean-basic} (\texttt{Pseudocontexts/Basic.lean})
  defines coordinate rank-one operators, structure operators, the
  root-of-unity phase construction, same-ray predicates, the revised
  plane-genuineness certificate, and trace identities.
\end{itemize}

\newpage
\tableofcontents

\newpage
\section{Project Configuration Files}

\subsubsection*{\texttt{lean-toolchain}}
\begin{lstlisting}[language=lean, numbers=none]
leanprover/lean4:v4.30.0
\end{lstlisting}

\subsubsection*{\texttt{lakefile.lean}}
\lstinputlisting{C:/mytex/lean/2026-qip26/lakefile.lean}

\subsubsection*{\texttt{Pseudocontexts.lean}}
\lstinputlisting{C:/mytex/lean/2026-qip26/Pseudocontexts.lean}

\subsubsection*{\texttt{verify.ps1}}
\lstinputlisting[numbers=none]{C:/mytex/lean/2026-qip26/verify.ps1}

\newpage
\section{Formalized Pseudocontext Algebra}
\label{app:lean-basic}

\subsubsection*{\texttt{Pseudocontexts/Basic.lean}}
\lstinputlisting{C:/mytex/lean/2026-qip26/Pseudocontexts/Basic.lean}

\end{document}
